\section{Vấn đề và động lực}\label{sec:introduction}

Trợ lý lập trình AI dạng tác tử (\emph{coding agent}) như Claude Code, Codex hay Antigravity tự đọc
mã, sửa file và chạy test trong kho mã của người dùng. Cách giao việc phổ biến nhất là \emph{vibe
coding}: mô tả việc cần làm trong vài câu rồi để tác tử viết mã ngay \cite{chou_2025,fawzy_2026}. Đề ngắn
như vậy bỏ ngỏ nhiều chi tiết, và tác tử phải tự đoán. Phát triển theo đặc tả hứa bù cho chỗ đó: tác tử
viết tài liệu thiết kế và kế hoạch trước, viết test trước khi viết mã, rồi kiểm lại cả bộ test
\cite{agarwal_2026,diaz_2026}; bộ skill Superpowers mà An dùng khi làm ViPharmind là một quy trình như vậy.

Bằng chứng định lượng cho lời hứa này còn ít. Nghiên cứu về vibe coding chủ yếu là định tính
\cite{chou_2025,fawzy_2026}; nghiên cứu về đặc tả dựa trên phỏng vấn \cite{agarwal_2026}, chấm bằng mô hình
ngôn ngữ \cite{taghavi_2026}, hoặc quan sát pull request, và nghiên cứu quan sát lớn nhất không thấy pull
request có đặc tả ít lỗi hơn \cite{hill_2026}. Các bộ đo như SWE-bench và FeatBench so sánh mô hình với
nhau, không so sánh quy trình làm việc \cite{jimenez_2024,chen_2025}. Chưa có thí nghiệm nào giữ nguyên đề
bài, công cụ và mô hình, chỉ đổi lời dặn về quy trình, rồi chấm bằng test mà tác tử không thấy.

Đề tài làm thí nghiệm đó trên ViPharmind, trợ lý AI tra cứu thuốc cho người Việt do Doãn Văn An xây dựng với
trợ lý lập trình AI ở mọi khâu (\NumCommits{} commit). Đề tài có hai phần: phần hồi cứu phân tích lịch sử git
của dự án, và phần thực nghiệm giao cùng một nhiệm vụ cho tác tử theo hai cách làm. Báo cáo lần 1 tập trung
vào phần thực nghiệm; phần hồi cứu và tổng quan tài liệu đầy đủ ở báo cáo cuối kỳ.
